\section{Introduction}
The $L(h,k)$-labeling problem assigns nonnegative integers to graph
vertices, with prescribed separation at distances one and two. It
originates in frequency assignment, where nearby transmitters require
different levels of interference protection \cite{griggs,calamoneri}.
The objective is to minimize the span between the largest and smallest
labels. Formulas are known for particular graph classes and parameter
regimes; a general-purpose exact method must also handle inputs for which
these formulas do not determine the optimum.

This study develops and evaluates an exact SAT pipeline for finite simple
undirected graphs with $h\geq k\geq1$. Boolean threshold variables encode
labels, and a bounded span search maintains a feasible solution together
with a lower bound. Assignment integer linear programming (ILP) provides
the principal exact baseline. Both methods can prove optimality, but may
retain an open interval when their time budget expires. We therefore
compare feasible-solution quality and proof coverage alongside runtime.
``General graphs'' describes the formulation's input domain, not a claim
of efficient solution on every graph.

Four research questions organize the study:
\begin{enumerate}
\item\label{rq:correctness} Does the SAT encoding preserve feasibility,
and does the span search report optimality only when it is established?
\item\label{rq:theory} Where exact formulas or proven bounds apply, do
computed labelings agree with them, and which instances can be certified
by matching lower and upper bounds?
\item\label{rq:comparison} How do runtime, model size and optimality
coverage compare between SAT and assignment ILP across the evaluated
graph sizes and families?
\item\label{rq:symmetry} For families with justified symmetries, how do
symmetry-breaking constraints affect CNF size and solution time?
\end{enumerate}

The study makes four contributions:
\begin{enumerate}
\item A general-graph SAT pipeline with encoding equivalence, a
span-search invariant and justified symmetry conditions.
\item Independent labeling checks and reference-bound comparisons,
with retained graphs, configurations and witnesses for reproducibility.
\item A comparison with assignment ILP on 117 fixed graph/parameter
instances that reports unresolved cases, bound quality and proof
coverage alongside runtime. A separate paired $L(2,1)$ study evaluates
symmetry effects by family under its original backend and budget.
\item A supplementary certificate pipeline combining mathematical
bounds, explicit labelings and independently replayed UNSAT proofs.
Clique and color-order certificates first close five of the 21 intervals
left open by pooled solver evidence. Weighted capacity bounds and a
color-rank-residual SAT formulation then close the remaining 16, yielding
independently checked optimal spans for all 117 instances.
\end{enumerate}

SAT and ILP formulations for distance-constrained labeling already exist
\cite{shao2013distance,shao2014strong}, and order encoding is established
\cite{tamura2009order}. Our contribution is the verified formulation and
the accompanying evidence on exact search, bound quality and symmetry
effects under stated protocols. The computational certificates establish
optima for the evaluated instances; they do not establish new formulas
over entire graph families.
